\documentclass[11pt,a4paper]{article}

%====================================================
% PACKAGES
%====================================================
\usepackage[utf8]{inputenc}
\usepackage[T1]{fontenc}
\usepackage[french]{babel}
\usepackage[a4paper,margin=2cm]{geometry}
\usepackage{amsmath,amssymb,mathtools,amsthm}
\usepackage{physics}
\usepackage{tikz}
\usepackage{fancyhdr}
\usepackage{lmodern}
\usepackage{enumitem}

\pagestyle{fancy}
\fancyhf{}
\fancyhead[L]{Cinématique du point}
\fancyhead[C]{Changement de bases}
\fancyhead[R]{\thepage}

\begin{document}

\begin{center}
\Huge \textbf{CINÉMATIQUE DU POINT MATÉRIEL}\\
\Large Version complète + démonstrations + algèbre linéaire
\end{center}

%====================================================
\section{Introduction complète}
%====================================================

La mécanique classique constitue un cadre fondamental de la physique permettant
de décrire le mouvement des corps macroscopiques.

Dans ce chapitre, on se limite à l’étude de la \textbf{cinématique},
c’est-à-dire à la description des mouvements indépendamment de leurs causes :contentReference[oaicite:1]{index=1}.

\bigskip

\textbf{Idée centrale :}
\begin{center}
Décrire un mouvement revient à exprimer un vecteur position dans une base adaptée.
\end{center}

\bigskip

Cela conduit à deux notions fondamentales :

\begin{itemize}
\item le choix d’un \textbf{référentiel}
\item le choix d’un \textbf{système de coordonnées}
\end{itemize}

\bigskip

\textbf{Remarque fondamentale :}

\begin{center}
Un vecteur est invariant, mais ses coordonnées dépendent du choix de la base.
\end{center}

---

%====================================================
\section{Référentiel et repère}
%====================================================

\subsection{Définition d’un référentiel}

Un référentiel est un ensemble de points matériels considérés comme fixes entre eux.

\bigskip

Il est nécessaire de préciser le référentiel pour définir une vitesse.

\bigskip

\textbf{Exemple :}

Dans un train en mouvement :
\begin{itemize}
\item vitesse par rapport au train
\item vitesse par rapport au sol
\end{itemize}

Ces deux vitesses sont différentes.

---

\subsection{Différence repère / référentiel}

\begin{itemize}
\item référentiel = notion physique
\item repère = outil mathématique
\end{itemize}

\bigskip

Un même référentiel peut être associé à différents repères.

---

%====================================================
\section{Repérage d’un point}
%====================================================

On choisit :
\begin{itemize}
\item un point origine $O$
\item une base $(\vec{e}_1,\vec{e}_2,\vec{e}_3)$
\end{itemize}

\[
\vec{OM} = u_1 \vec{e}_1 + u_2 \vec{e}_2 + u_3 \vec{e}_3
\]

---

\subsection{Base orthonormée directe}

Une base est orthonormée si :

\[
\vec{e}_i \cdot \vec{e}_j = 0 \quad (i \neq j)
\]

\[
||\vec{e}_i|| = 1
\]

Elle est directe si :

\[
\vec{e}_1 \wedge \vec{e}_2 = \vec{e}_3
\]

---

%====================================================
\section{Coordonnées cartésiennes}
%====================================================

\[
\vec{OM} = x\vec{e}_x + y\vec{e}_y + z\vec{e}_z
\]

\subsection{Schéma détaillé}

\begin{center}
\begin{tikzpicture}[scale=2]

\draw[->] (0,0) -- (2,0) node[right]{$x$};
\draw[->] (0,0) -- (0,2) node[above]{$y$};
\draw[->] (0,0) -- (-1,-1) node[left]{$z$};

\draw[->] (0,0) -- (1.2,1.3) node[right]{$M$};

\end{tikzpicture}
\end{center}

---

%====================================================
\section{Coordonnées cylindriques}
%====================================================

\subsection{Définition physique}

Utilisé pour les mouvements avec symétrie de rotation.

\[
x = r\cos\theta, \quad y = r\sin\theta
\]

\[
\vec{OM} = r\vec{e}_r + z\vec{e}_z
\]

---

\subsection{Schéma complet}

\begin{center}
\begin{tikzpicture}[scale=2]

\draw (0,0) circle (1.5);

\draw[->] (0,0) -- (1.2,0.9);
\node at (1.3,1) {$M$};

\draw[->] (0,0) -- (2,0) node[right]{$x$};
\draw[->] (0,0) -- (0,2) node[above]{$y$};

\draw (0.5,0) arc (0:35:0.5);
\node at (0.7,0.3){$\theta$};

\node at (1.2,0.6){$r$};

\end{tikzpicture}
\end{center}

---

%====================================================
\section{Base cylindrique — démonstration}
%====================================================

\subsection{Objectif}

Exprimer $(\vec{e}_r,\vec{e}_\theta)$ en fonction de $(\vec{e}_x,\vec{e}_y)$.

---

\subsection{Démonstration géométrique}

On projette :

\[
\vec{e}_r = \cos\theta \vec{e}_x + \sin\theta \vec{e}_y
\]

---

\[
\vec{e}_\theta = -\sin\theta \vec{e}_x + \cos\theta \vec{e}_y
\]

---

\subsection{Matrice associée}

\[
P =
\begin{pmatrix}
\cos\theta & -\sin\theta \\
\sin\theta & \cos\theta
\end{pmatrix}
\]

---

\subsection{Interprétation}

Cette matrice est une rotation.

---

%====================================================
\section{Coordonnées sphériques}
%====================================================

\[
x = r\sin\theta\cos\varphi
\]

\[
y = r\sin\theta\sin\varphi
\]

\[
z = r\cos\theta
\]

---

\subsection{Schéma}

\begin{center}
\begin{tikzpicture}[scale=2]

\draw (0,0) circle (1.2);
\draw[->] (0,0) -- (0.7,0.9);
\node at (0.8,0.9){$r$};

\end{tikzpicture}
\end{center}

---

%====================================================
\section{Conclusion partie 1}
%====================================================

On a introduit :

\begin{itemize}
\item les systèmes de coordonnées
\item les bases associées
\item le changement de base géométrique
\end{itemize}

\bigskip

La suite consiste à étudier :

\begin{itemize}
\item la dérivation des vecteurs
\item la vitesse
\item l’accélération
\end{itemize}

\section{Dérivation d’un vecteur dans un référentiel}

\subsection{Problème fondamental}

On souhaite calculer la dérivée temporelle d’un vecteur :

\[
\vec{v}(t) = \sum_i a_i(t)\vec{e}_i(t)
\]

\bigskip

\textbf{Problème :}

\begin{itemize}
\item les coefficients $a_i(t)$ dépendent du temps
\item MAIS aussi les vecteurs $\vec{e}_i(t)$
\end{itemize}

\bigskip

Donc :

\[
\frac{d\vec{v}}{dt} \neq \sum_i \dot{a}_i \vec{e}_i
\]

---

\subsection{Formule générale}

\[
\boxed{
\frac{d\vec{v}}{dt}
=
\sum_i \dot{a}_i \vec{e}_i + \sum_i a_i \frac{d\vec{e}_i}{dt}
}
\]

\bigskip

C’est LA formule clé du chapitre.

---

%====================================================
\section{Dérivation de la base cylindrique}
%====================================================

\subsection{Expression de $\vec{e}_r$}

\[
\vec{e}_r = \cos\theta \vec{e}_x + \sin\theta \vec{e}_y
\]

---

\subsection{Dérivation complète (ligne par ligne)}

On dérive :

\[
\frac{d\vec{e}_r}{dt}
=
\frac{d}{dt}(\cos\theta)\vec{e}_x + \frac{d}{dt}(\sin\theta)\vec{e}_y
\]

Or :

\[
\frac{d}{dt}(\cos\theta) = -\sin\theta \dot{\theta}
\]

\[
\frac{d}{dt}(\sin\theta) = \cos\theta \dot{\theta}
\]

Donc :

\[
\frac{d\vec{e}_r}{dt}
=
-\sin\theta \dot{\theta} \vec{e}_x + \cos\theta \dot{\theta} \vec{e}_y
\]

Factorisation :

\[
\frac{d\vec{e}_r}{dt}
=
\dot{\theta}(-\sin\theta \vec{e}_x + \cos\theta \vec{e}_y)
\]

Or :

\[
-\sin\theta \vec{e}_x + \cos\theta \vec{e}_y = \vec{e}_\theta
\]

Donc :

\[
\boxed{
\frac{d\vec{e}_r}{dt} = \dot{\theta} \vec{e}_\theta
}
\]

---

\subsection{Dérivation de $\vec{e}_\theta$}

\[
\vec{e}_\theta = -\sin\theta \vec{e}_x + \cos\theta \vec{e}_y
\]

On dérive :

\[
\frac{d\vec{e}_\theta}{dt}
=
-\cos\theta \dot{\theta} \vec{e}_x - \sin\theta \dot{\theta} \vec{e}_y
\]

Factorisation :

\[
\frac{d\vec{e}_\theta}{dt}
=
-\dot{\theta}(\cos\theta \vec{e}_x + \sin\theta \vec{e}_y)
\]

Or :

\[
\cos\theta \vec{e}_x + \sin\theta \vec{e}_y = \vec{e}_r
\]

Donc :

\[
\boxed{
\frac{d\vec{e}_\theta}{dt} = -\dot{\theta} \vec{e}_r
}
\]

---

%====================================================
\section{Interprétation géométrique}
%====================================================

\begin{center}
\begin{tikzpicture}[scale=2]

\draw[->] (0,0) -- (1,0) node[right]{$\vec{e}_r$};
\draw[->] (0,0) -- (0,1) node[above]{$\vec{e}_\theta$};

\draw[->,red] (1,0) -- (1,0.4) node[right]{$d\vec{e}_r$};

\end{tikzpicture}
\end{center}

\bigskip

\textbf{Conclusion :}

\begin{itemize}
\item $\vec{e}_r$ tourne
\item sa dérivée est perpendiculaire
\end{itemize}

\bigskip

 C’est une rotation instantanée.

---

%====================================================
\section{Vitesse d’un point}
%====================================================

\subsection{Définition}

\[
\vec{v} = \frac{d\vec{OM}}{dt}
\]

---

%====================================================
\subsection{Méthode 1 : calcul direct}
%====================================================

\[
\vec{OM} = r\vec{e}_r + z\vec{e}_z
\]

\[
\vec{v} = \frac{d}{dt}(r\vec{e}_r + z\vec{e}_z)
\]

\[
\vec{v} = \dot{r}\vec{e}_r + r\frac{d\vec{e}_r}{dt} + \dot{z}\vec{e}_z
\]

Remplacement :

\[
\vec{v} = \dot{r}\vec{e}_r + r\dot{\theta}\vec{e}_\theta + \dot{z}\vec{e}_z
\]

---

%====================================================
\subsection{Méthode 2 : méthode géométrique}
%====================================================

On considère un déplacement infinitésimal :

\[
d\vec{OM} = dr\vec{e}_r + r d\theta \vec{e}_\theta + dz\vec{e}_z
\]

\bigskip

 Interprétation :

\begin{itemize}
\item variation de $r$ → déplacement radial
\item variation de $\theta$ → déplacement circulaire
\item variation de $z$ → déplacement vertical
\end{itemize}

---

En divisant par $dt$ :

\[
\vec{v} =
\dot{r}\vec{e}_r + r\dot{\theta}\vec{e}_\theta + \dot{z}\vec{e}_z
\]

---

%====================================================
\section{Interprétation physique de la vitesse}
%====================================================

\[
\vec{v} =
\dot{r}\vec{e}_r + r\dot{\theta}\vec{e}_\theta + \dot{z}\vec{e}_z
\]

\bigskip

\begin{itemize}
\item $\dot{r}$ : vitesse radiale
\item $r\dot{\theta}$ : vitesse tangentielle
\item $\dot{z}$ : vitesse verticale
\end{itemize}

---

%====================================================
\section{Lien avec rotation}
%====================================================

On peut écrire :

\[
\frac{d\vec{e}_r}{dt} = \vec{\Omega} \wedge \vec{e}_r
\]

avec :

\[
\vec{\Omega} = \dot{\theta} \vec{e}_z
\]

---

\textbf{Interprétation :}

\begin{center}
La base cylindrique tourne avec une vitesse angulaire $\dot{\theta}$.
\end{center}

---

%====================================================
\section{Conclusion de la partie}
%====================================================

On a montré :

\begin{itemize}
\item comment dériver une base mobile
\item comment obtenir la vitesse par deux méthodes
\item le lien avec la rotation
\end{itemize}

\bigskip

 La prochaine étape est :

\begin{center}
\textbf{le calcul de l’accélération (beaucoup plus riche)}
\end{center}

%====================================================
\section{Repère de Frenet}
%====================================================

\subsection{Motivation}

Dans les systèmes précédents (cartésien, cylindrique, sphérique),
la base dépend de la position dans l’espace.

\bigskip

Le repère de Frenet adopte une approche différente :

\begin{center}
\textbf{la base est adaptée à la trajectoire elle-même}
\end{center}

---

%====================================================
\subsection{Définition du repère de Frenet}
%====================================================

On considère une trajectoire paramétrée par le temps $t$ :

\[
\vec{OM}(t)
\]

On définit :

---

\textbf{Vecteur tangent :}

\[
\vec{T} = \frac{\vec{v}}{||\vec{v}||}
\]

---

\textbf{Vecteur normal :}

\[
\vec{N} = \frac{1}{||\dot{\vec{T}}||} \dot{\vec{T}}
\]

---

\textbf{Vecteur binormal :}

\[
\vec{B} = \vec{T} \wedge \vec{N}
\]

---

Le triplet :

\[
(\vec{T}, \vec{N}, \vec{B})
\]

forme une base orthonormée directe.

---

%====================================================
\subsection{Schéma}
%====================================================

\begin{center}
\begin{tikzpicture}[scale=2]

\draw (0,0) .. controls (0.5,0.8) .. (1,0);

\draw[->] (0.5,0.4) -- (0.9,0.4) node[right] {$\vec{T}$};
\draw[->] (0.5,0.4) -- (0.5,0.8) node[above] {$\vec{N}$};

\end{tikzpicture}
\end{center}

---

%====================================================
\subsection{Formules de Frenet}
%====================================================

On montre que :

\[
\boxed{
\frac{d\vec{T}}{dt} = \kappa v \vec{N}
}
\]

\[
\boxed{
\frac{d\vec{N}}{dt} = -\kappa v \vec{T} + \tau v \vec{B}
}
\]

\[
\boxed{
\frac{d\vec{B}}{dt} = -\tau v \vec{N}
}
\]

---

\subsection{Signification des grandeurs}

\begin{itemize}
\item $\kappa$ : courbure
\item $\tau$ : torsion
\item $v = ||\vec{v}||$ : vitesse
\end{itemize}

---

%====================================================
\subsection{Démonstration de la première formule}
%====================================================

On a :

\[
\vec{T} = \frac{\vec{v}}{||\vec{v}||}
\]

On dérive :

\[
\frac{d\vec{T}}{dt}
=
\frac{\dot{\vec{v}}}{||\vec{v}||}
-
\frac{\vec{v}}{||\vec{v}||^2} \frac{d||\vec{v}||}{dt}
\]

Après simplification, on montre que :

\[
\frac{d\vec{T}}{dt} \perp \vec{T}
\]

Donc :

\[
\frac{d\vec{T}}{dt} = \kappa v \vec{N}
\]

---

%====================================================
\subsection{Accélération dans le repère de Frenet}
%====================================================

\[
\vec{v} = v \vec{T}
\]

On dérive :

\[
\vec{a} = \dot{v}\vec{T} + v \frac{d\vec{T}}{dt}
\]

\[
\vec{a} = \dot{v}\vec{T} + \kappa v^2 \vec{N}
\]

---

\subsection{Interprétation physique}

\begin{itemize}
\item $\dot{v}\vec{T}$ : accélération tangentielle
\item $\kappa v^2 \vec{N}$ : accélération normale (centripète)
\end{itemize}

---

%====================================================
\subsection{Lien avec le mouvement circulaire}
%====================================================

Pour un cercle de rayon $R$ :

\[
\kappa = \frac{1}{R}
\]

Donc :

\[
\vec{a} = \frac{v^2}{R} \vec{N}
\]

�� On retrouve l’accélération centripète.

---

%====================================================
\section{Conclusion sur le repère de Frenet}
%====================================================

\begin{itemize}
\item il est adapté à la trajectoire
\item il simplifie l’expression de l’accélération
\item il relie directement géométrie et mécanique
\end{itemize}

\bigskip

\begin{center}
\textbf{C’est le repère le plus naturel pour décrire un mouvement.}
\end{center}

%====================================================
\section{Accélération d’un point}
%====================================================

\subsection{Définition}

L’accélération est définie comme la dérivée de la vitesse :

\[
\vec{a} = \frac{d\vec{v}}{dt}
\]

\bigskip

Contrairement à la vitesse, il n’existe pas de méthode géométrique simple :
il faut dériver explicitement.

---

%====================================================
\section{Accélération en coordonnées cylindriques}
%====================================================

\subsection{Rappel de la vitesse}

\[
\vec{v} =
\dot{r}\vec{e}_r + r\dot{\theta}\vec{e}_\theta + \dot{z}\vec{e}_z
\]

---

\subsection{Dérivation complète}

On dérive chaque terme :

---

\textbf{Premier terme :}

\[
\frac{d}{dt}(\dot{r}\vec{e}_r)
=
\ddot{r}\vec{e}_r + \dot{r}\frac{d\vec{e}_r}{dt}
\]

\[
= \ddot{r}\vec{e}_r + \dot{r}\dot{\theta}\vec{e}_\theta
\]

---

\textbf{Deuxième terme :}

\[
\frac{d}{dt}(r\dot{\theta}\vec{e}_\theta)
\]

On applique la règle du produit :

\[
= \frac{d}{dt}(r\dot{\theta})\vec{e}_\theta
+ r\dot{\theta}\frac{d\vec{e}_\theta}{dt}
\]

Or :

\[
\frac{d}{dt}(r\dot{\theta}) = \dot{r}\dot{\theta} + r\ddot{\theta}
\]

et :

\[
\frac{d\vec{e}_\theta}{dt} = -\dot{\theta}\vec{e}_r
\]

Donc :

\[
= (\dot{r}\dot{\theta} + r\ddot{\theta})\vec{e}_\theta
- r\dot{\theta}^2 \vec{e}_r
\]

---

\textbf{Troisième terme :}

\[
\frac{d}{dt}(\dot{z}\vec{e}_z) = \ddot{z}\vec{e}_z
\]

(car $\vec{e}_z$ est constant)

---

%====================================================
\subsection{Regroupement final}
%====================================================

On regroupe les composantes :

---

\textbf{Composante radiale :}

\[
\ddot{r} - r\dot{\theta}^2
\]

---

\textbf{Composante orthoradiale :}

\[
2\dot{r}\dot{\theta} + r\ddot{\theta}
\]

---

\textbf{Composante verticale :}

\[
\ddot{z}
\]

---

\[
\boxed{
\vec{a} =
(\ddot{r} - r\dot{\theta}^2)\vec{e}_r
+
(2\dot{r}\dot{\theta} + r\ddot{\theta})\vec{e}_\theta
+
\ddot{z}\vec{e}_z
}
\]

---

%====================================================
\section{Interprétation physique}
%====================================================

\subsection{Terme centripète}

\[
-r\dot{\theta}^2
\]

\begin{itemize}
\item dirigé vers le centre
\item responsable du mouvement circulaire
\end{itemize}

---

\subsection{Terme de Coriolis}

\[
2\dot{r}\dot{\theta}
\]

\begin{itemize}
\item apparaît si le rayon varie
\item lié au changement de base
\end{itemize}

---

\subsection{Terme angulaire}

\[
r\ddot{\theta}
\]

accélération de rotation.

---

%====================================================
\section{Coordonnées sphériques}
%====================================================

\subsection{Rappel du système}

\[
(r,\theta,\varphi)
\]

\begin{itemize}
\item $r$ : distance à l’origine
\item $\theta$ : colatitude
\item $\varphi$ : azimut
\end{itemize}

---

\subsection{Vitesse (rappel)}

\[
\vec{v} =
\dot{r}\vec{e}_r
+
r\dot{\theta}\vec{e}_\theta
+
r\sin\theta \dot{\varphi}\vec{e}_\varphi
\]

---

%====================================================
\section{Démonstration de l’accélération sphérique}
%====================================================

 Calcul long mais fondamental.

---

On dérive :

\[
\vec{a} = \frac{d\vec{v}}{dt}
\]

---

\textbf{1er terme :}

\[
\frac{d}{dt}(\dot{r}\vec{e}_r)
=
\ddot{r}\vec{e}_r + \dot{r}\frac{d\vec{e}_r}{dt}
\]

---

\textbf{2e terme :}

\[
\frac{d}{dt}(r\dot{\theta}\vec{e}_\theta)
=
(\dot{r}\dot{\theta} + r\ddot{\theta})\vec{e}_\theta
+ r\dot{\theta}\frac{d\vec{e}_\theta}{dt}
\]

---

\textbf{3e terme :}

\[
\frac{d}{dt}(r\sin\theta \dot{\varphi}\vec{e}_\varphi)
\]

développement :

\[
=
(\dot{r}\sin\theta \dot{\varphi}
+ r\cos\theta \dot{\theta}\dot{\varphi}
+ r\sin\theta \ddot{\varphi})\vec{e}_\varphi
+ r\sin\theta \dot{\varphi}\frac{d\vec{e}_\varphi}{dt}
\]

---

%====================================================
\subsection{Résultat final}
%====================================================

\[
\boxed{
\vec{a} =
(\ddot{r} - r\dot{\theta}^2 - r\sin^2\theta \dot{\varphi}^2)\vec{e}_r
}
\]

\[
+
\boxed{
(2\dot{r}\dot{\theta} + r\ddot{\theta}
- r\sin\theta\cos\theta \dot{\varphi}^2)\vec{e}_\theta
}
\]

\[
+
\boxed{
(2\dot{r}\sin\theta \dot{\varphi}
+ 2r\cos\theta\dot{\theta}\dot{\varphi}
+ r\sin\theta\ddot{\varphi})\vec{e}_\varphi
}
\]

---

%====================================================
\section{Interprétation physique (sphérique)}
%====================================================

\begin{itemize}
\item accélération radiale → attraction vers centre
\item termes mixtes → couplage rotation
\item termes en $\sin\theta$ → géométrie sphérique
\end{itemize}

---

%====================================================
\section{Conclusion de la partie}
%====================================================

On a obtenu :

\begin{itemize}
\item accélération cylindrique complète
\item accélération sphérique complète
\item interprétation physique des termes
\end{itemize}

\bigskip

 La dernière étape :

\begin{center}
\textbf{relier tout cela à l’algèbre linéaire (matrices)}
\end{center}

%====================================================
\section{Changement de base : point de vue algèbre linéaire}
%====================================================

\subsection{Cadre général}

Soit $V$ un espace vectoriel de dimension 3.

On considère deux bases :

\[
\mathcal{B} = (\vec{e}_1,\vec{e}_2,\vec{e}_3), \quad
\mathcal{B}' = (\vec{e}_1',\vec{e}_2',\vec{e}_3')
\]

---

\subsection{Définition de la matrice de passage}

La matrice de passage de $\mathcal{B}$ vers $\mathcal{B}'$ est définie par :

\[
P_{\mathcal{B},\mathcal{B}'}
=
\begin{pmatrix}
| & | & | \\
[\vec{e}_1']_{\mathcal{B}} &
[\vec{e}_2']_{\mathcal{B}} &
[\vec{e}_3']_{\mathcal{B}} \\
| & | & |
\end{pmatrix}
\]

---

\subsection{Interprétation fondamentale}

\[
P_{\mathcal{B},\mathcal{B}'} = \text{Mat}_{\mathcal{B}',\mathcal{B}}(\mathrm{Id}_V)
\]

 C’est la matrice de l’identité entre deux bases.

---

%====================================================
\section{Transformation des coordonnées}
%====================================================

Soit un vecteur $\vec{v}$.

\[
[\vec{v}]_{\mathcal{B}} = P [\vec{v}]_{\mathcal{B}'}
\]

\[
[\vec{v}]_{\mathcal{B}'} = P^{-1} [\vec{v}]_{\mathcal{B}}
\]

---

\subsection{Démonstration complète}

On écrit :

\[
\vec{v} = \sum_i v_i' \vec{e}_i'
\]

Or chaque $\vec{e}_i'$ s’exprime dans $\mathcal{B}$ :

\[
\vec{e}_i' = \sum_j P_{ji} \vec{e}_j
\]

Donc :

\[
\vec{v} = \sum_i v_i' \sum_j P_{ji} \vec{e}_j
\]

\[
= \sum_j \left( \sum_i P_{ji} v_i' \right)\vec{e}_j
\]

Donc :

\[
[\vec{v}]_{\mathcal{B}} = P [\vec{v}]_{\mathcal{B}'}
\]

---

%====================================================
\section{Changement de matrice d’un endomorphisme}
%====================================================

Soit $f : V \to V$.

On note :

\[
A = \text{Mat}_{\mathcal{B}}(f), \quad
B = \text{Mat}_{\mathcal{B}'}(f)
\]

---

\subsection{Résultat fondamental}

\[
\boxed{
B = P^{-1} A P
}
\]

---

\subsection{Démonstration}

On a :

\[
f(\vec{v}) = A [\vec{v}]_{\mathcal{B}}
\]

Mais :

\[
[\vec{v}]_{\mathcal{B}} = P [\vec{v}]_{\mathcal{B}'}
\]

Donc :

\[
f(\vec{v}) = A P [\vec{v}]_{\mathcal{B}'}
\]

Or :

\[
[f(\vec{v})]_{\mathcal{B}'} = P^{-1} A P [\vec{v}]_{\mathcal{B}'}
\]

Donc :

\[
B = P^{-1} A P
\]

---

%====================================================
\section{Matrices semblables}
%====================================================

\textbf{Définition :}

Deux matrices $A$ et $B$ sont semblables si :

\[
B = P^{-1} A P
\]

---

\textbf{Interprétation :}

\begin{center}
Même transformation, bases différentes.
\end{center}

---

%====================================================
\section{Lien avec la mécanique}
%====================================================

En mécanique :

\[
\vec{v}_{cart} = P(t) \vec{v}_{mobile}
\]

 La matrice dépend du temps !

---

%====================================================
\section{Dérivation matricielle}
%====================================================

On dérive :

\[
\frac{d\vec{v}}{dt}
=
\frac{d}{dt}(P(t)\vec{v}')
\]

---

\subsection{Règle du produit}

\[
= \dot{P}(t)\vec{v}' + P(t)\dot{\vec{v}}'
\]

---

\subsection{Réécriture}

On veut exprimer en coordonnées physiques :

\[
\vec{v} = P\vec{v}'
\]

Donc :

\[
\vec{v}' = P^{-1}\vec{v}
\]

---

\subsection{Substitution}

\[
\frac{d\vec{v}}{dt}
=
P\dot{\vec{v}}' + \dot{P}P^{-1}\vec{v}
\]

---

%====================================================
\section{Définition de l’opérateur de rotation}
%====================================================

On pose :

\[
\boxed{
\Omega = \dot{P}P^{-1}
}
\]

---

\subsection{Formule fondamentale}

\[
\boxed{
\frac{d\vec{v}}{dt}
=
\frac{d\vec{v}}{dt}_{mobile}
+
\Omega \vec{v}
}
\]

---

%====================================================
\section{Interprétation physique}
%====================================================

\[
\Omega \vec{v} = \vec{\Omega} \wedge \vec{v}
\]

---

\subsection{Conclusion}

\begin{itemize}
\item $\Omega$ représente la rotation de la base
\item c’est une vitesse angulaire
\end{itemize}

---

%====================================================
\section{Lien avec la base cylindrique}
%====================================================

\[
\vec{\Omega} = \dot{\theta} \vec{e}_z
\]

\[
\frac{d\vec{e}_r}{dt} = \vec{\Omega} \wedge \vec{e}_r
\]

---

%====================================================
\section{Conclusion générale}
%====================================================

\begin{center}
\Huge \textbf{Synthèse}
\end{center}

\bigskip

\begin{itemize}
\item Le changement de base est une transformation linéaire
\item En mécanique, cette transformation dépend du temps
\item Cela crée des termes supplémentaires (Coriolis, centripète)
\item Ces termes ont une interprétation géométrique : rotation
\end{itemize}

\bigskip

\begin{center}
\Huge \textbf{La mécanique est une géométrie dynamique}
\end{center}

%====================================================
\section{Exercices d'application}
%====================================================

\subsection*{Exercice 1 — Changement de base cylindrique (*)}

On considère un point $M$ dans le plan $(Oxy)$ de coordonnées cartésiennes $(x,y)$.

\begin{enumerate}
\item Exprimer les coordonnées cylindriques $(r,\theta)$ en fonction de $(x,y)$.
\item Exprimer les vecteurs $\vec{e}_r$ et $\vec{e}_\theta$ en fonction de $\vec{e}_x$ et $\vec{e}_y$.
\item Écrire la matrice de passage $P$ de la base cartésienne vers la base cylindrique.
\item Montrer que cette matrice est orthogonale.
\end{enumerate}

\textbf{Objectif :} comprendre le lien entre rotation et changement de base.

---

%====================================================
\subsection*{Exercice 2 — Dérivation des vecteurs de base (**)}

On considère la base cylindrique $(\vec{e}_r,\vec{e}_\theta)$.

\begin{enumerate}
\item Repartir de l’expression de $\vec{e}_r$ dans la base cartésienne.
\item Calculer explicitement $\frac{d\vec{e}_r}{dt}$.
\item Montrer que cette dérivée est orthogonale à $\vec{e}_r$.
\item Interpréter géométriquement ce résultat.
\end{enumerate}

\textbf{Objectif :} comprendre la notion de base mobile.

---

%====================================================
\subsection*{Exercice 3 — Vitesse en coordonnées cylindriques (**)}

On considère un point dont la position est donnée par :

\[
\vec{OM} = r(t)\vec{e}_r
\]

\begin{enumerate}
\item Calculer la vitesse par dérivation directe.
\item Retrouver le résultat par la méthode géométrique.
\item Identifier les composantes radiale et tangentielle.
\end{enumerate}

\textbf{Objectif :} maîtriser les deux méthodes.

---

%====================================================
\subsection*{Exercice 4 — Mouvement circulaire uniforme (***)}

On considère un mouvement circulaire de rayon $R$ avec :

\[
r = R, \quad \theta = \omega t
\]

\begin{enumerate}
\item Calculer la vitesse.
\item Calculer l’accélération.
\item Montrer que l’accélération est dirigée vers le centre.
\item Interpréter physiquement le résultat.
\end{enumerate}

\textbf{Objectif :} comprendre l’accélération centripète.

---

%====================================================
\subsection*{Exercice 5 — Accélération complète (***)}

On considère un mouvement général :

\[
r = r(t), \quad \theta = \theta(t)
\]

\begin{enumerate}
\item Retrouver l’expression complète de l’accélération.
\item Identifier les termes :
\begin{itemize}
\item radial
\item orthoradial
\end{itemize}
\item Interpréter le terme $2\dot{r}\dot{\theta}$.
\end{enumerate}

\textbf{Objectif :} comprendre le terme de Coriolis.

---

%====================================================
\subsection*{Exercice 6 — Coordonnées sphériques (****)}

On considère un point en coordonnées sphériques :

\[
(r,\theta,\varphi)
\]

\begin{enumerate}
\item Écrire le vecteur position.
\item Déterminer la vitesse par la méthode géométrique.
\item Identifier les trois composantes de la vitesse.
\item Interpréter physiquement chaque terme.
\end{enumerate}

\textbf{Objectif :} comprendre la géométrie sphérique.

---

%====================================================
\subsection*{Exercice 7 — Matrice de passage (****)}

Soit deux bases :

\[
\mathcal{B} = (\vec{e}_x,\vec{e}_y), \quad
\mathcal{B}' = (\vec{e}_r,\vec{e}_\theta)
\]

\begin{enumerate}
\item Écrire la matrice de passage $P$.
\item Calculer $P^{-1}$.
\item Vérifier que $P^{-1} = P^T$.
\item Interpréter ce résultat.
\end{enumerate}

\textbf{Objectif :} lien rotation ↔ matrices orthogonales.

---

%====================================================
\subsection*{Exercice 8 — Endomorphisme et changement de base (****)}

Soit une transformation linéaire de matrice :

\[
A =
\begin{pmatrix}
1 & 2 \\
0 & 1
\end{pmatrix}
\]

\begin{enumerate}
\item Calculer la matrice $B = P^{-1}AP$ où $P$ est une rotation.
\item Interpréter le résultat.
\item Montrer que $A$ et $B$ sont semblables.
\end{enumerate}

\textbf{Objectif :} lien avec l’algèbre linéaire.

---

%====================================================
\subsection*{Exercice 9 — Rotation et dérivation (*****)}

On considère une matrice de rotation $R(t)$.

\begin{enumerate}
\item Montrer que $R^T R = I$.
\item Dériver cette relation.
\item En déduire que $\dot{R}R^{-1}$ est antisymétrique.
\end{enumerate}

\textbf{Objectif :} comprendre la structure de $\Omega$.

---

%====================================================
\subsection*{Exercice 10 — Formulation générale (*****)}

On considère :

\[
\vec{v} = R(t)\vec{v}'
\]

\begin{enumerate}
\item Dériver cette expression.
\item Introduire $\Omega = \dot{R}R^{-1}$.
\item Montrer que :

\[
\frac{d\vec{v}}{dt} =
\frac{d\vec{v}}{dt}_{mobile} + \Omega \vec{v}
\]

\item Interpréter physiquement cette formule.
\end{enumerate}

\textbf{Objectif :} niveau ENS.

---

%====================================================
\section{Conclusion des exercices}
%====================================================

Ces exercices permettent de :

\begin{itemize}
\item maîtriser les changements de base
\item comprendre la dérivation dans une base mobile
\item relier mécanique et algèbre linéaire
\end{itemize}

\bigskip

\begin{center}
\textbf{La compréhension vient autant des calculs que de l’interprétation.}
\end{center}
\end{document}